\section*{\centering\large{HALAMAN PENGESAHAN}}
\vspace{0.08in}
\begin{center}
\textbf{\large{RADIASI HAWKING PADA RUANG-WAKTU LUBANG HITAM VAIDYA-KERR}}\\
\vspace{0.08in}
Disusun untuk memenuhi salah satu syarat \\
memperoleh gelar Sarjana pada\\ 
Program Studi S-1 Fisika\\
Departemen Fisika \\
Fakultas Sains dan Analitika Data\\
Institut Teknologi Sepuluh Nopember \\
\vspace{15mm}
 Oleh \\
\vspace{1mm}
\textbf{MUHAMMAD HANIF RAHMANI} \\
\vspace{1mm}
\textbf{01111940000074} \\

\vspace{0.3in}
Surabaya, 1 April 2023
\end{center}
\vspace{0.3in}
%\raggedleft Surabya, 1 Desember 2017 \\
\begin{center}
	\textbf{Disetujui oleh Pembimbing Proposal Tugas Akhir:} \\
	\vspace{2mm}
 \end{center}
	\begin{center}
\vspace{2mm}
	    Dosen Pembimbing \\
	    Departemen Fisika FSAD ITS \\
	    \vspace{0.8in}
	    \textbf{\underline{Dr.rer.nat Bintoro A Subagyo}} \\	
	    \textbf{NIP.19790719 200501.1.015} \\
\end{center}

